\section{Introduction}
\label{sec:intro}

Double Higgs boson production is the process through which the shape of the Higgs
potential is probed directly, and it is by now approaching the sensitivity needed to
observe the Standard Model (SM) rate~\cite{ATLAS:2024ish,CMS_combination}. As that
sensitivity is reached, the characterization of the signal becomes as important as its
observation: deviations of the Higgs self-coupling alone are a rare possibility in
nature, since new physics in the Higgs sector generally modifies several couplings at
once~\cite{Durieux:2022hbu}, and each modification changes both the rate and the topology
of the process.

Among the effective field theory descriptions of these deviations, the Standard Model
Effective Field Theory (SMEFT) is the more restrictive one, in that the \HH and
single-Higgs couplings to fermions and gluons are correlated rather than independent. Its
parameter space is nonetheless large, and it becomes larger once the subleading operators
--- those entering only through loop-induced diagrams --- are included, as they must be
for a consistent treatment at NLO QCD~\cite{Heinrich:2023rsd}. Simulating every point of
that space, and processing every point through an experimental analysis, is not feasible.

Two complementary techniques make the problem tractable, and this report applies both to
\ggHH production in SMEFT with leading and subleading operators. The first is event-level
reweighting (ELR), which parameterizes the differential cross section in the Wilson
coefficients and transports a reference sample to any point of the parameter space,
removing the need to simulate each point; we follow the technique established for the
same process in Ref.~\cite{ELR_R5}. The second is the matrix method, which expresses the
signal at any point as a linear combination of a fixed, minimal set of templates,
removing the need to process each point through the analysis chain; here we follow the
technique of Ref.~\cite{HH_nonggf_R5} rather than the one of Ref.~\cite{HHMatrix_R5}, the
two differing in how the template basis is searched and scored, for the reasons given in
Sec.~\ref{sec:matrix}.

What the two give together is the subject of this report. The parameterization is solved
on a basis of forty templates, every one of which lies in a plane in which a single pair
of coefficients is varied, so that each template belongs to a plane the results are
already presented in (Sec.~\ref{sec:matrix}). Three results follow from combining the two
techniques on that basis, and each reduces what an analysis has to simulate, process or
derive for itself.

The first is the cost of the basis. Forty templates do not require forty simulated
samples: each template can be built from a few generated samples by reweighting, combined
bin by bin. Five samples are enough. Doing so reduces the median residual against the
generated samples from $3.04\%$ to $2.03\%$ and the worst case from $70\%$ to $5\%$, and
reaching the same precision by generating would take about eight hundred samples
(Sec.~\ref{sec:validation_nsamples}).

The second is the collision energy. A change of $\sqrt{s}$ acts on the initial state
alone, so between nearby energies a grid can be transported with parton densities instead
of being generated again. Measured against the points generated at both $13$ and
$13.6$~TeV, the transport reproduces the spectrum to $0.2\%$ per bin below $1$~TeV once a
single normalization constant is absorbed, and the $14$~TeV coefficients provided here
are obtained that way. The step to $100$~TeV is not of that kind --- the spectrum extends
to much higher \mHH and the relative weight of the diagrams shifts with it --- and the
coefficients there are extracted from a grid generated at that energy
(Sec.~\ref{sec:energy_transport}).

The third is the uncertainty. The interpolation uncertainty of a morphed prediction
separates into weights, which belong to the basis and are published once, and template
errors, which belong to whichever analysis applies them. An analysis therefore obtains
its own band from its own templates with nothing about the effective field theory
re-derived; applied to two selections of the same \HHbbbb release the same weights give
bands differing by a factor five to ten (Sec.~\ref{sec:validation_realistic_band}). The
prescription is given with the covariance of the templates, which is the form a basis
built by reweighting one set of events requires (Sec.~\ref{sec:uncertainties}).

Three prescriptions complete what an \HH analysis needs in order to interpret its own
result in the Wilson coefficients, without importing conventions from a single-Higgs
combination. The order of the expansion is applied once to $\sigma\times\mathrm{BR}$ and
not to its factors, which is a question of consistency rather than of convention and
costs up to a factor of two at the edges of the coefficient ranges
(Sec.~\ref{sec:trunc_product}). The range over which a chosen dimension-6 truncation is
one an analysis would want is given per operator and per \mHH bin
(Sec.~\ref{sec:trunc_convergence}). And $c_{\PH\mathrm{kin}}$ has a physical region:
the linearly truncated partial widths stay positive only for
$c_{\PH\mathrm{kin}}>-\Lambda^2/2v^2$, which is $-8.25$ at $\Lambda=1$~TeV and $-2.06$
for the \HH triangle term, and an \HH analysis can constrain that coefficient with data
it already has by giving its \ttH-enriched regions the same coefficient dependence as
the signal (Sec.~\ref{sec:trunc_singleh}).

The structure of the parameterization is also what allows a fit in one or two
coefficients to drop most of the templates and run faster, at an unchanged result
(Sec.~\ref{sec:slice_models}).

Results are provided at four centre-of-mass energies. The $\sqrt{s}=13.6$~TeV results are
the main results of this report, matching the ongoing LHC Run 3 programme. Results at
$\sqrt{s}=13$~TeV are also provided, since a large number of analyses are still
reinterpreting the LHC Run 2 dataset; they are so similar to the $13.6$~TeV ones that only
the latter are shown in the figures, the $13$~TeV numbers being given with the ancillary
material. Coefficients at $\sqrt{s}=14$~TeV are obtained by transport rather than from a
grid of their own. Finally, results at $\sqrt{s}=100$~TeV are provided to enable
feasibility studies for a future circular hadron collider such as
FCC-hh~\cite{FCC:2018vvp}. \textcolor{red}{TODO: check the FCC-hh reference} The $13$ and
$13.6$~TeV extractions share the same grid of generated points, the $13$~TeV one being a
subset of the $13.6$~TeV one; the $100$~TeV extraction uses a grid of its own, with
narrower coefficient ranges and coarser scans (Sec.~\ref{sec:grids}).

Only the elements needed to make this report self-contained are repeated here; for the
derivation and validation of the two techniques we refer to Refs.~\cite{ELR_R5}
and~\cite{HH_nonggf_R5}.

This report is organized as follows. Section~\ref{sec:framework} introduces the SMEFT
operators considered and the diagrams they affect, and Sec.~\ref{sec:parameterization}
gives the parameterization of the differential cross section and the reweighting formula.
Section~\ref{sec:eventsimulation} describes the simulation setup and the generated grids,
and Sec.~\ref{sec:extraction} the extraction of the $A$-coefficients.
Section~\ref{sec:matrix} presents the matrix method, the restriction of the basis to the
scan planes and the basis optimization itself. Section~\ref{sec:validation} validates
both methods on the generated samples with no selection applied, and
Sec.~\ref{sec:validation_realistic} under the selections of a public \HHbbbb analysis,
where the number of templates a fit has to carry is also reduced.
Section~\ref{sec:energy_transport} transports the parameterization to another collision
energy and Sec.~\ref{sec:uncertainties} gives the uncertainty prescription. The last
three sections concern the expansion rather than the machinery:
Sec.~\ref{sec:trunc_modelling} shows that both routes reproduce a truncated prediction,
Sec.~\ref{sec:truncation} bounds where the dimension-6 interpretation holds and how it
meets a single-Higgs combination, and Sec.~\ref{sec:rge} states how the coefficients are
run between scales. Section~\ref{sec:conclusion} summarizes.
